% das Papierformat zuerst
%
%\documentclass[a4paper, 11pt,bibtotoc]{scrartcl}
\documentclass[a4paper, 11pt,bibtotoc,abstracton]{scrreprt}  

%Für URL
\usepackage{url}
\renewcommand{\UrlFont}{\rmfamily}

%Für Zitaten
\usepackage{cite}

%Abkürzungen
\usepackage{acronym}


%Inhaltsverzeichnis bearbeiten
\usepackage{tocbibind}

% für mathematische Symbole
\usepackage{amsmath}


% Vektorgrafiken mit Latex importieren
\usepackage{import}

\usepackage[colorlinks=false, pdfborder={0 0 0}]{hyperref}
%documentclass[a4paper, 12pt]{article}  
% deutsche Silbentrennung
\usepackage[english,ngerman]{babel}
\renewcommand{\sectfont}{\rmfamily\bfseries}
% wegen deutschen Umlauten
\usepackage[utf8]{inputenc}
\usepackage{graphicx}
\usepackage{subfigure}\hyphenation{Bit-rate}

%Algorithm schreiben
\usepackage{algorithm2e}

%Tabellen
\usepackage{booktabs}
\usepackage{multirow}
\usepackage{colortbl}


%Farben
\usepackage{color}
\usepackage{listings}%Code einbinden
\definecolor{darkblue}{rgb}{.08,.21,.36}
\definecolor{darkred}{rgb}{.6,.19,.20}
\definecolor{darkgreen}{rgb}{0,.6,0}
\definecolor{red}{rgb}{.98,0,0}
\definecolor{lightblue}{rgb}{0.8,0.85,1}
%\definecolor{lightgrey}{rgb}{0.98,0.98,0.98}
\definecolor{lightgrey}{gray}{.98}
\definecolor{black}{rgb}{0.0,0.0,0.0}


\lstloadlanguages{C++}
\lstset{%
  language=C++,
  basicstyle=\small,
  commentstyle=\itshape\color{darkgreen},
  keywordstyle=\bfseries\color{darkblue},
  stringstyle=\color{darkred},
  showspaces=false,
  showtabs=false,
  columns=fixed,
  backgroundcolor=\color{lightgrey},
  numbers=left,
  frame=single,
  numberstyle=\tiny,
  breaklines=true,
  showstringspaces=false,
  xleftmargin=1cm,
  basicstyle=\small
}%

\usepackage{amssymb}%Mathematische Symbole, wie R,N,Q,Z,...
\setlength{\parindent}{0pt} %einrücken nach absatz verhindern
%\usepackage{setspace}%Zeilenabstand
\usepackage{algorithmic}%Für Pseudocode

%%%%%%%%%%%%%%%%%%%%%%%%%%%%%%%%%%%
%Seiten Kopf- und Fußzeilen
\usepackage[automark,						
		headsepline,								
		plainfootsepline, 
		]{scrlayer-scrpage}

\automark[section]{chapter} 
\pagestyle{scrheadings}			

\clearscrheadings	%Alte Kopfformatierungen entfernen
\clearscrplain		%Alte Plain-Formatierung entfernen
\clearscrheadfoot %Alten Fuß entfernen
\cfoot[\pagemark]{\pagemark}%Seitenzahl zentriert im Fuß 
\ihead{\leftmark}
\ohead{\rightmark} 
 
%%%%%%%%%%%%%%%%%%%%%%%%%%%%%%%%%%%

\begin{document}
\begin{titlepage}
\begin{center}
{\huge \textbf{Philipps-Universität Marburg}}\\[0.5cm]
\textbf{Fachbereich 12 - Mathematik und Informatik}\\[0.5cm]

\begin{figure}[h]
	\centering
		\includegraphics[width=0.8\textwidth]{fig/unilogo.pdf}
\end{figure}

{\huge \textbf{{\large \\[1cm]Diplomarbeit}}}
\\[1cm]

{\Huge \textbf{Progressive Meshes\\ mit CUDA}}
\\[1cm]

{\large von}\\
{\large Max Mustermann}\\
{\large Dezember 2009}\\[3cm]

{\large
Betreuer:\\ Prof. Dr. Thorsten Thormählen\\[1cm]
Arbeitsgruppe Grafik und Multimedia Programmierung}

\end{center}
\end{titlepage}
\newpage
\thispagestyle{empty}
\section*{}
\newpage
\thispagestyle{empty}
\vspace*{7cm}
\textbf{\Large {Erklärung}}\\[0.5cm]
Ich, Max Mustermann (Informatikstudent an der Philipps-Universität Marburg, Matrikelnummer
xxxxxx), versichere an Eides statt, dass ich die vorliegende Diplomarbeit selbstständig verfasst und keine anderen
als die angegebenen Quellen und Hilfsmittel verwendet habe. Die hier vorliegende Diplomarbeit wurde weder in ihrer jetzigen noch in einer ähnlichen Form einer Prüfungskommission vorgelegt.\\[1cm]
Marburg, 17. Dezember 2009\\[0.5cm]
Max Mustermann
\newpage
\shipout\null

%\setstretch {1.15}%Zeilenabstand setzen

  
% hier beginnt das Dokument